\documentclass[10pt,a4paper]{amsart}
\usepackage[margin=19mm]{geometry}
\usepackage{amsmath,amssymb,mathtools}
\usepackage[scr=rsfs,cal=euler]{mathalfa}
\usepackage{xfrac}
\usepackage[unicode,hidelinks,bookmarks=false]{hyperref}
\newcommand{\ie}{i.e.\ }
\newcommand{\cf}{cf.\ }
\newcommand{\resp}{respectively\ }
\newcommand{\Subsection}{\subsection}
\providecommand{\nobreakdash}{\nobreak\hbox{-}}
\setlength{\emergencystretch}{2em}
\multlinegap0pt
\usepackage{amssymb}
\usepackage[shortlabels]{enumitem}
\usepackage{dsfont}
\usepackage{xcolor}
\usepackage{csquotes}
\usepackage{float}
\usepackage{tikz}
\newtheorem{lem}{Lemme}
\title{Action de groupe algébrique sur la compactification hybride}
\author{Alexandre Roy}
\date{Source: arXiv:2512.00201v2 (CC BY 4.0)}
\begin{document}
\maketitle

\noindent\emph{Source and licence.} Action de groupe algébrique sur la compactification hybride. Version arXiv:2512.00201v2 (CC BY 4.0), distributed by arXiv under the Creative Commons Attribution 4.0 International licence, http://creativecommons.org/licenses/by/4.0/, as recorded in the arXiv licence metadata for this exact version and shown on the abstract page https://arxiv.org/abs/2512.00201v2. Full licence notice: \texttt{licenses/arxiv-2512.00201-CC-BY-4.0.txt}.\par\smallskip

\noindent\textit{Editorial scope.} Counted unit: the article's lem lem (source0.tex lines 1227--1229). The article's complete original proof of it is delivered with it (source0.tex lines 1230--1246, one \texttt{proof} environment of 17 lines). No mathematics is written, repaired or re-derived: the statement and the proof are the article's own lines, in the article's own notation, and no retained line is substituted or re-flowed.  No further source-local context is needed: every cross-reference of every retained line resolves inside the two delivered spans.  Nothing was omitted from any retained span, so the package carries no omission record.  No retained line contains a citation, so the wrapper carries no bibliography and no bibliography entry.  No retained line contains a figure environment, an image-inclusion command or drawing code, so no image is delivered and the package declares no asset dependency.  Editorial formatting only, in addition: an AMS \texttt{amsart} preamble that restates the article's own declarations and macros used by the retained lines, namely the environments \texttt{lem} on separate counters, together with the standard packages those lines need.  The retained environments print in this document as Lemma 1 in order of appearance, whereas the article prints the same items with the numbering of its own full document.  The complete native source of the article as distributed in the arXiv e-print for this exact version is bundled unchanged as \texttt{sources/190032/source0.tex}.
\begin{lem}
Il existe $\beta$ algébrique sur $\mathcal{H}(x)$ dont le polynôme minimal est à coefficients dans $Frac(R/ker~\eta_x)$ avec $y \in (\mathcal{H}(x)(\beta))^m$.
\end{lem}

\begin{proof}[Démonstration du lemme]
On sait que $y \in (\mathcal{H}(x)(y_{1}, \cdots, y_{m}))^m$, donc par le théorème de l'élément primitif, il existe $\alpha$ algébrique sur $\mathcal{H}(x)$ tel que $\mathcal{H}(x)(y_{1}, \cdots, y_{m}) = \mathcal{H}(x)(\alpha)$.

Soit $\epsilon > 0$, notons $\mu_\alpha = \sum_{k=0}^l \mu_{\alpha,k}T^k$ le polynôme minimal de $\alpha$ et $\alpha_i$ ses racines. Alors il existe $P = \sum_{k=0}^l P_k T^K \in Frac(R/ker~\eta_x)[T]$ tel que $max(\tilde{\eta_x}(P_k - \mu_{\alpha,k})) \leq \epsilon$ où $\tilde{\eta_x}$ est la norme sur $K$ qui prolonge $\eta_x$. Si l'on prends $\mu_\alpha$ unitaire, on peut aussi prendre $P$ unitaire et ainsi si $\beta$ est une racine de $P$, on a :

\begin{equation*}
\tilde{\eta_x}(\beta) \leq max(\eta_x(P_0), \cdots, \eta_x(P_{l-1}), 1) \leq  max(\tilde{\eta_x}(\mu_{\alpha,0}), \cdots, \tilde{\eta_x}(\mu_{\alpha,l-1}), 1) +\epsilon :=C.
\end{equation*}

D'où
\begin{equation*}
C^l\epsilon \geq \tilde{\eta_x}((P-\mu_{\alpha})(\beta)) = \tilde{\eta_x}(\mu_{\alpha}(\beta)) = \prod_{\alpha_i~racines~de~\mu_\alpha} \tilde{\eta_x}(\alpha_i-\beta). 
\end{equation*}
Donc, l'un des $\alpha_i$ vérifie que $\tilde{\eta_x}(\alpha_i-\beta) \leq C \epsilon^{\frac{1}{l}}$. En prenant $\epsilon$ tel que $\epsilon < (\frac{max (\tilde{\eta_x}(\alpha_i -\alpha_j))}{C})^l$, on obtient un unique $\alpha_i$ vérifiant cette condition et telle que $\tilde{\eta_x}(\alpha_i-\beta) < min_{i\neq j}(\tilde{\eta_x}(\alpha_j-\beta))$.

Donc par le lemme de Krasner, on a $\mathcal{H}(x)(\alpha) \subset \mathcal{H}(x)(\beta)$. Comme le polynôme $P$ est annulateur de $\beta$, on a $deg(\mu_\beta) \leq deg(P_\beta) = deg(\mu_\alpha)$, où $\mu_\beta$ désigne un polynôme minimal de $\beta$. Par l'inclusion ci-dessus, on a  $deg(\mu_\beta) \geq deg(\mu_\alpha)$ et donc $P$ est en fait un polynôme minimal de $\beta$. 
\end{proof}

\end{document}
